\chapter{Before you start}

\begin{prereq}
This book assumes the following, and never re-teaches it.
\begin{itemize}
  \item \textbf{Calculus:} derivatives, the chain rule, integrals, integration by parts, substitution, Taylor series, $e^x=\lim_{n\to\infty}(1+x/n)^n$.
  \item \textbf{Series:} geometric sums $1+q+\dots+q^{n-1}=\frac{1-q^n}{1-q}$ for $q\ne1$, and convergent series.
  \item \textbf{Sets:} union, intersection, complement, $A\subseteq B$, the empty set $\varnothing$.
  \item \textbf{Matrices:} products, transposes, determinants, solving a linear system.
\end{itemize}
\end{prereq}

\section*{How to read this book}
\begin{itemize}
  \item \textbf{Revise from the cards.} Every chapter ends with a formula sheet of cards. A card holds one general formula, boxed. Each numbered row below it is one situation: what you are given, what to plug in, the result, when it is valid, and one worked line with numbers. Every smaller formula in the chapters is a row of some card, and no formula appears on two cards. The index below lists every card; reading all of them in order is a full revision of the course.
  \item \textbf{To plug in:} find the card whose box covers the question, pick the row that matches what you are given, check its ``valid when'' line, and substitute.
  \item \textbf{To skip a chapter,} read only its formula sheet: the chapter's first box gives the page. A chapter with no card of its own names the card that covers it.
  \item \textbf{Problems are in the final Part,} grouped by topic. Every solution names its card and row (a link), plugs in the numbers, and reads off the answer.
  \item \textbf{Part 0 is the probability review.} If you have revised it, skip it. Every later chapter lists, in its first box, what it uses from Part 0, so you can jump back to one item.
  \item \textbf{Proofs are optional reading.} Each theorem is stated first, with what it means. Its proof sits just below, as numbered steps. Nothing later needs you to have read a proof. A few proofs need mathematics beyond this book; those are left out, and the book says so.
  \item \textbf{Every term is defined, in bold, before it is used,} and the glossary at the end lists every defined term with its page.
  \item \textbf{Every expectation names the probability it uses:} $\E_P$ under the real-world probability $P$, and $\E_{P^*}$ under the risk-neutral probability $P^*$ of Part~4.
\end{itemize}

\tcblistof[\section*]{cards}{Index of cards}
